\section{Setup, observables, and how regimes are named}
\label{sec:setup}

\paragraph{Protocol.}
Unless stated otherwise: $N=40$ agents in two equal classes, $K=30$, agenda
$P=5$ (simple, $\alpha=P/K=0.17$) or $P=100$ (complex, $\alpha=3.3$), the
symmetric field \eqref{eq:disc-field}, discriminators drawn independently of
class. Initial state $\hat{w}\sim\mathcal{N}(0,I)$, $C=I$, $V=1$ and
$\mu\sim\mathcal{U}(-1,1)$, the last chosen so that a fresh population starts
with as many frustrated triples as unfrustrated ones. Emitter, receiver and
issue are drawn uniformly at each step, and every society is run for $500$
interactions per ordered pair. $K$ is recovered rather than chosen and the
interaction count is calibrated; Appendix~\ref{app:simulation} gives the
provenance of every number, along with the sensitivity of each observable to the
measurement time --- which, since the dynamics anneals rather than converging
(\S\ref{sec:model}), is a genuine parameter and not a convergence tolerance.

\textbf{Every reported quantity is a mean over $24$ independent replicates per
grid point}, each with its own initial condition, agenda, class assignment,
discriminator draw and interaction schedule, on a $49\times49$ grid in
$(d,\fd)$. Uncertainties are replicate standard errors for per-point quantities.
For anything derived from the map as a whole --- a boundary location, a regime
area, a crossover --- we resample whole replicate sheets rather than propagating
per-point errors, because replicates within a sheet share an interaction
schedule and their errors are therefore correlated
(Appendix~\ref{app:simulation}).

\paragraph{Observables.}
Two pairwise quantities carry everything. The \emph{ideological overlap}
$\qov_{IJ}=\cos(w_I,w_J)$ is $+1$ when two agents agree on every issue and $-1$
when they disagree on every one. The \emph{directed trust}
$\tdir=1-2\Phi(\hmu)$ is $+1$ for a fully trusted source and $-1$ for a fully
distrusted one. With the signed class relation $s_{IJ}=\kappa_I\kappa_J$ and
$\bar{t}_{IJ}=(t_{I|J}+t_{J|I})/2$,
%
\begin{equation}
  \CcT = \bigl\langle s_{IJ}\,\bar{t}_{IJ} \bigr\rangle,
  \qquad
  \CcO = \bigl\langle s_{IJ}\,\qov_{IJ} \bigr\rangle,
  \qquad
  \CtO = \bigl\langle \bar{t}_{IJ}\,\qov_{IJ} \bigr\rangle,
  \label{eq:correlations}
\end{equation}
%
averaged over unordered pairs. These are a \emph{renaming} of the three
correlations used in earlier presentations of this model, not a correction: the
quantities are numerically identical, and the accompanying code asserts equality
to $10^{-12}$. We rename them only because $\CcT$, $\CcO$, $\CtO$ say which
sectors are being correlated, and the earlier symbols do not.

\paragraph{The class correlations are not sufficient, and this matters for one
whole region of the diagram.}
A population can be sharply polarised along an axis unrelated to class, in which
case all of \eqref{eq:correlations} are near zero --- and a class-only analysis
would report it as structureless. We therefore also measure class-independent
polarisation. Let $\lambda_a$ be the eigenvalues of the centred overlap matrix
and set
%
\begin{equation}
  P \;=\; \lambda_1^{2}\big/\textstyle\sum_a \lambda_a^{2},
  \qquad
  \Po \;=\; \frac{P-P_{\text{null}}}{P_{\max}(N)-P_{\text{null}}},
  \label{eq:polarisation}
\end{equation}
%
and identically $\Pt$ on the symmetrised trust matrix. $P$ is largest for a
two-bloc (rank-one) structure and about $1/\mathrm{rank}$ for an isotropic one.
The squared-eigenvalue fraction is used rather than a variance fraction because
the centred trust matrix is indefinite, where a variance fraction is
meaningless; using one formula for both sectors is what makes the comparison in
\S\ref{sec:agenda} a comparison of like with like. Both the null $P_{\text{null}}$,
measured at initialisation, and the attainable maximum
$P_{\max}(N)=(N-1)^2/\bigl((N-1)^2+N-2\bigr)$ are $N$-dependent, so both are
divided out; without that, the finite-size panels would show a drift in
polarisation with $N$ that is purely an artefact of the estimator.

We also measure social balance over triples
\citep{heider1946attitudes,cartwright1956structural}, which separates
\emph{organised} disagreement from disorder: two coherent, mutually opposed blocs
give balance $1$ whatever they are made of, and a population that cannot settle
goes negative. We report the continuous scores $\Bo,\Bt$ and their sign-based
counterparts separately, and only ever describe the latter as a fraction of
balanced triples --- a continuous product of overlaps is not one. Both are
computed in closed form from matrix products rather than by enumerating triples
(Appendix~\ref{app:observables}).

\paragraph{Regimes are classified, not labelled by hand.}
Rather than placing names on a picture, we fix thresholds, apply them to each
replicate separately, and report the modal regime together with the fraction of
replicates that agree. Thresholds are calibrated from the no-discrimination
null: the cut on $|\CcT|$ is three standard deviations of $\CcT$ where $d=0$, so
``class-correlated'' is measured against how large the correlation gets when
there is provably no class signal, at the same $N$. Classifying replicates
rather than the replicate mean matters --- a point where half the replicates fall
in one regime and half in another is bistable, and averaging first would render
it as an intermediate value and hide that. Points where fewer than $70\%$ of
replicates agree are drawn hatched, which is the boundary uncertainty made
visible. Appendix~\ref{app:observables} gives the thresholds, the rule, and the
regime areas recomputed at $\{2/3,1,3/2\}$ times each cut.

\paragraph{On naming.}
Regime names describe measured behaviour. In particular we do not use
``spin glass'': that name requires overlap distributions, metastability and
initialisation dependence, none of which a pair correlation establishes, so the
region with $d<0$ is named for what is actually measured --- negative affective
balance with no ideological alignment. We likewise write ``regimes'' and
``crossovers'' rather than ``phases'' and ``transitions'' except where the
finite-size evidence of \S\ref{sec:regimes} supports the stronger word.
